\documentclass[conference]{IEEEtran}
\IEEEoverridecommandlockouts
% The preceding line is only needed to identify funding in the first footnote. If that is unneeded, please comment it out.
%Template version as of 6/27/2024

\usepackage{cite}
\usepackage{amsmath,amssymb,amsfonts}
\usepackage{algorithmic}
\usepackage{graphicx}
\usepackage{textcomp}
\usepackage{xcolor}
\def\BibTeX{{\rm B\kern-.05em{\sc i\kern-.025em b}\kern-.08em
    T\kern-.1667em\lower.7ex\hbox{E}\kern-.125emX}}
\begin{document}

\title{Conference Paper Title*\\
{\footnotesize \textsuperscript{*}Note: Sub-titles are not captured for https://ieeexplore.ieee.org  and
should not be used}
\thanks{Identify applicable funding agency here. If none, delete this.}
}

\author{\IEEEauthorblockN{1\textsuperscript{st} Given Name Surname}
\IEEEauthorblockA{\textit{dept. name of organization (of Aff.)} \\
\textit{name of organization (of Aff.)}\\
City, Country \\
email address or ORCID}
\and
\IEEEauthorblockN{2\textsuperscript{nd} Given Name Surname}
\IEEEauthorblockA{\textit{dept. name of organization (of Aff.)} \\
\textit{name of organization (of Aff.)}\\
City, Country \\
email address or ORCID}
\and
\IEEEauthorblockN{3\textsuperscript{rd} Given Name Surname}
\IEEEauthorblockA{\textit{dept. name of organization (of Aff.)} \\
\textit{name of organization (of Aff.)}\\
City, Country \\
email address or ORCID}
\and
\IEEEauthorblockN{4\textsuperscript{th} Given Name Surname}
\IEEEauthorblockA{\textit{dept. name of organization (of Aff.)} \\
\textit{name of organization (of Aff.)}\\
City, Country \\
email address or ORCID}
\and
\IEEEauthorblockN{5\textsuperscript{th} Given Name Surname}
\IEEEauthorblockA{\textit{dept. name of organization (of Aff.)} \\
\textit{name of organization (of Aff.)}\\
City, Country \\
email address or ORCID}
\and
\IEEEauthorblockN{6\textsuperscript{th} Given Name Surname}
\IEEEauthorblockA{\textit{dept. name of organization (of Aff.)} \\
\textit{name of organization (of Aff.)}\\
City, Country \\
email address or ORCID}
}

\maketitle

\begin{abstract}
This document is a model and instructions for \LaTeX.
This and the IEEEtran.cls file define the components of your paper [title, text, heads, etc.]. *CRITICAL: Do Not Use Symbols, Special Characters, Footnotes, 
or Math in Paper Title or Abstract.
\end{abstract}

\begin{IEEEkeywords}
component, formatting, style, styling, insert.
\end{IEEEkeywords}

\section{Introduction}
Automatic time complexity analysis is a fundamental problem in software engineering, program verification, and online judging systems. Accurate complexity classification is crucial for detecting inefficient implementations, preventing time limit exceeded (TLE) errors, and supporting intelligent grading and code review systems. Traditionally, complexity analysis relies on static program analysis techniques such as abstract interpretation, control-flow graph (CFG) inspection, and symbolic reasoning. Although these methods provide sound upper bounds, they are often overly conservative, difficult to scale to real-world programs, and expensive to implement for modern programming languages. 

Recently, large language models (LLMs) have demonstrated remarkable capabilities in code understanding and reasoning. Several studies have explored using LLMs to infer algorithmic complexity directly from source code. However, despite their strong semantic understanding, LLM-based complexity prediction still suffers from severe hallucination and systematic bias. In particular, we observe that LLMs frequently underestimate high-order polynomial complexities (e.g., $O(n^3)$ and $O(n^4)$), often confusing them with quadratic or near-linear patterns. Such misjudgments significantly limit the reliability of purely LLM-driven approaches. 

An alternative direction is dynamic complexity estimation, which infers asymptotic growth by executing programs under scaled input sizes and fitting runtime measurements to predefined growth models. While dynamic fitting can capture empirical growth trends, it is inherently sensitive to runtime noise, limited input ranges, and short-scale behaviors. In practice, dynamic fitting alone may produce unstable or misleading classifications, especially when the input size cannot reach the asymptotic regime due to time or memory constraints. 

These limitations motivate a hybrid approach that integrates static semantic reasoning and dynamic empirical measurement. In this paper, we propose a unified framework that combines LLM-based static complexity prediction with dynamic runtime curve fitting. By leveraging complementary strengths of both paradigms, our framework corrects systematic LLM biases using runtime evidence while stabilizing dynamic fitting results with semantic priors derived from LLM reasoning. 

We conduct extensive experiments on a dataset of 4,900 competitive programming solutions collected from Codeforces. Our results show that LLM-only static prediction achieves an overall accuracy of 63.7\%, but exhibits pronounced underestimation on cubic and higher-order polynomial programs. By introducing our fusion strategy, we significantly improve classification accuracy, particularly on high-order complexity classes. 

In summary, this paper makes the following contributions: 

\begin{itemize} 
\item We present a hybrid complexity classification framework that integrates LLM-based static reasoning with dynamic runtime curve fitting. 
\item We construct a large-scale complexity-labeled dataset of 4,900 competitive programming solutions for empirical evaluation. 
\item We reveal systematic underestimation biases of LLMs on high-order polynomial complexities and demonstrate that our fusion strategy effectively mitigates these errors. 
\end{itemize} 

This work provides a practical and scalable direction toward robust, automatic time complexity classification and establishes a foundation for future research on LLM-assisted program performance analysis.

\section{Related Work} 
 \label{sec:related_work} 
 
 \subsection{Benchmarks for Time Complexity Prediction} 
 Early studies on program time-complexity prediction mainly relied on relatively small or restricted benchmarks. 
 For example, CoRCoD contains 929 annotated \emph{Java} programs and augments them with hand-engineered features 
 (e.g., counts of loops, methods, variables, and specific data-structure usage) for machine learning-based prediction. :contentReference[oaicite:0]{index=0} 
 However, prior benchmarks are limited in scale and language coverage, and often do not fully capture the role of 
 \emph{input representation} (e.g., distinguishing numeric constants from input size indicators), which is crucial for 
 worst-case time complexity reasoning. 
 
 To address these limitations, Baik et al.\ introduced \textsc{CodeComplex}, a large-scale benchmark designed 
 specifically for evaluating LLM reasoning on worst-case time complexity prediction. :contentReference[oaicite:1]{index=1} 
 \textsc{CodeComplex} highlights several advantages over earlier benchmarks: it is larger and more diverse (covering 
 multiple languages rather than only Java), includes a more comprehensive labeled complexity taxonomy, explicitly 
 models input representation, and proposes evaluation metrics beyond plain class accuracy for nuanced analysis. :contentReference[oaicite:2]{index=2} 
 In our work, we adopt \textsc{CodeComplex} as the primary benchmark and construct a modified subset tailored to our 
 execution-based setting (Section~\ref{sec:dataset}), focusing on a six-class label space 
 (\texttt{constant}, \texttt{logn}, \texttt{linear}, \texttt{nlogn}, \texttt{quadratic}, \texttt{cubic}). 
 
 \subsection{Learning- and LLM-based Complexity Prediction} 
 A common paradigm for complexity prediction is to treat it as a supervised classification problem from code text 
 (or code-derived features). Beyond traditional ML baselines (e.g., decision trees, random forests, SVMs), 
 recent work evaluates a wide range of pretrained programming language models (PLMs), including encoder-based models 
 (CodeBERT, GraphCodeBERT, UniXcoder) and encoder--decoder models (PLBART, CodeT5/CodeT5+), as well as both closed- 
 and open-source LLMs. :contentReference[oaicite:3]{index=3} 
 These studies show that even strong models may exhibit systematic errors on higher-order classes, motivating methods 
 that go beyond purely static text-based inference. 
 
 In parallel, the broader LLM reasoning literature explores prompting strategies such as zero-shot/few-shot prompting 
 and chain-of-thought prompting to elicit more structured inference. :contentReference[oaicite:4]{index=4} 
 While these techniques can improve reasoning in some settings, complexity prediction remains challenging due to the 
 need to jointly consider algorithmic structure, input size variables, and control-flow behavior. 
 
 \subsection{Evaluation Metrics for Hierarchical Complexity Labels} 
 Time complexity classes have a natural hierarchy: confusing $O(n\log n)$ with $O(n)$ is typically less severe than 
 confusing $O(n^3)$ with $O(1)$. To reflect this, \textsc{CodeComplex} proposes the \emph{Hierarchy Complexity Score} 
 (HC-Score), which penalizes wrong predictions proportionally to their hierarchical distance from the ground truth, 
 and a windowed variant (HC$_w$) that truncates penalties beyond a specified distance. :contentReference[oaicite:5]{index=5} 
 Following this line of work, we report standard class-based metrics (accuracy, macro-F1) and optionally include 
 hierarchy-aware metrics (HC/HC$_w$) to better characterize near-miss behaviors, especially around adjacent polynomial 
 classes. 
 
 \subsection{Static--Dynamic Hybrid Analysis} 
 Most existing benchmarks and baselines primarily study \emph{static} prediction from code (via features, PLMs, or LLMs). 
 In contrast, our work integrates LLM-based static reasoning with \emph{dynamic} runtime curve fitting, using execution 
 evidence to correct systematic static misclassifications and using semantic priors to stabilize fitting under noisy or 
 short-range measurements. This hybrid perspective complements prior benchmark-driven studies and targets robust 
 recognition of closely related polynomial classes such as \texttt{quadratic} vs.\ \texttt{cubic}.

\section{Methodology}
\label{sec:methodology}

\subsection{Fusion Strategy} 
 \label{subsec:fusion} 
 
 The LLM-based static prediction and the dynamic runtime fitting provide complementary signals for 
 time complexity inference. Static reasoning by LLMs captures high-level algorithmic intent but may 
 suffer from hallucination and systematic bias, while dynamic fitting reflects empirical growth trends 
 but is sensitive to runtime noise and limited input scales. To leverage the strengths of both, we adopt 
 an \emph{LLM-first with fitting-based correction} fusion strategy. 
 
 \subsubsection{Static Prior and Dynamic Evidence} 
 \label{subsubsec:prior_evidence} 
 
 We treat the LLM prediction as a \emph{semantic prior}: 
 \[ 
 \hat{y}_{\text{LLM}} = f_{\text{LLM}}(P), 
 \] 
 which reflects the model's reasoning over control flow, loop structure, and algorithmic patterns. 
 Dynamic runtime fitting provides \emph{empirical evidence}: 
 \[ 
 \hat{y}_{\text{fit}} = f_{\text{fit}}(\{(n_i, t_i)\}), 
 \] 
 where runtimes $t_i$ are measured using \texttt{time.perf\_counter()} over a sequence of input scales 
 $n_i \in \{3, 50, 100, 150, \dots, 3000\}$. 
 
 \subsubsection{Correction Rule} 
 \label{subsubsec:correction} 
 
 Our fusion rule prioritizes the static LLM prediction and applies dynamic fitting as a correction 
 signal only when there is strong empirical evidence of misclassification. Formally, the final prediction 
 $\hat{y}$ is determined as follows: 
 
 \begin{enumerate} 
     \item \textbf{Agreement:} If $\hat{y}_{\text{LLM}} = \hat{y}_{\text{fit}}$, we output 
     $\hat{y} = \hat{y}_{\text{LLM}}$. 
     \item \textbf{Dynamic failure:} If runtime fitting fails (e.g., insufficient valid measurements, 
     timeout, or unstable fitting), we output $\hat{y} = \hat{y}_{\text{LLM}}$. 
     \item \textbf{Conflict with strong evidence:} If $\hat{y}_{\text{LLM}} \neq \hat{y}_{\text{fit}}$, 
     we accept the fitted result $\hat{y}_{\text{fit}}$ only when the fitting confidence exceeds a 
     predefined threshold. Concretely, we require that the best-fitting model achieves a substantially 
     lower penalized error than alternative models on larger input scales. 
     \item \textbf{Fallback:} In all other cases, we retain $\hat{y} = \hat{y}_{\text{LLM}}$. 
 \end{enumerate} 
 
 \subsubsection{High-Order Bias Correction} 
 \label{subsubsec:high_order} 
 
 This correction mechanism is particularly effective for cubic-time programs. In our preliminary 
 analysis, LLMs exhibit a systematic tendency to underestimate cubic complexity as quadratic or 
 near-linear. When runtime measurements on larger input sizes (e.g., $n > 1000$) consistently favor 
 a cubic growth trend, the dynamic signal overrides the static prior, correcting this bias. Conversely, 
 for lower-order classes where runtime fitting is more susceptible to short-range noise, the LLM prior 
 provides a stabilizing effect. 
 
 Overall, this \emph{LLM-first with correction} strategy ensures robustness against both LLM 
 hallucination and overfitting in dynamic analysis, yielding more reliable complexity classification.

\section{Experiments}
\label{sec:experiments}

\subsection{Experimental Setup} 
 \label{subsec:exp_setup} 
 
 This subsection describes the evaluation protocol, runtime measurement procedure, input scaling 
 configuration, and model settings used in our experiments. 
 
 \subsubsection{Dataset and Label Space} 
 \label{subsubsec:setup_dataset} 
 All experiments are conducted on a modified version of the CodeComplex benchmark. 
 We focus on Python programs and a six-class taxonomy: 
 \texttt{constant}, \texttt{logn}, \texttt{linear}, \texttt{nlogn}, \texttt{quadratic}, and \texttt{cubic}. 
 We follow the benchmark's standard train/test split when available; otherwise, we create a split by 
 program ID to avoid leakage. Unless stated otherwise, all reported results are computed on the test set. 
 
 \subsubsection{Methods Compared} 
 \label{subsubsec:setup_methods} 
 We evaluate three approaches: 
 (i) \textbf{LLM-only}, which predicts a complexity class from source code via prompting; 
 (ii) \textbf{Fit-only}, which estimates the class by executing the program under scaled inputs and fitting 
 runtime curves; and (iii) \textbf{Fusion (ours)}, which follows an LLM-first with fitting-based correction 
 strategy (Section~\ref{subsec:fusion}). 
 
 \subsubsection{Runtime Measurement} 
 \label{subsubsec:setup_runtime} 
 For dynamic analysis, we execute each program under a sequence of input scales 
 $n \in \{3, 50, 100, 150, \dots, 3000\}$, where the sequence is monotonically increasing. 
 We measure runtime using \texttt{time.perf\_counter()} by taking the difference between the timestamps 
 immediately before and after the target function execution. To reduce the impact of system noise, we 
 repeat each run $r$ times and use the median runtime for each scale. 
 We enforce a per-run timeout $\tau$ and a per-program budget to prevent excessive execution time. 
 Programs that raise exceptions, time out, or produce invalid outputs at a scale are marked as failed for 
 that scale; programs with fewer than $k_{\min}$ valid scales are considered dynamic failures. 
 
 \subsubsection{Input Generation and Validity Checks} 
 \label{subsubsec:setup_inputs} 
 Because competitive-programming solutions may have heterogeneous input formats, we adopt a scalable 
 input generation strategy that produces syntactically valid inputs for each program and scale. 
 For each $n$, the generator instantiates the input template by substituting size-related parameters 
 with $n$ (or functions of $n$) while respecting common constraints (e.g., integer ranges and list lengths). 
 We apply lightweight validity checks (e.g., successful parsing and non-empty output) before recording a 
 measurement. When a program expects multiple test cases, we treat each test case independently and 
 aggregate runtimes according to the benchmark guideline. 
 
 \subsubsection{Curve Fitting Configuration} 
 \label{subsubsec:setup_fitting} 
 We fit measured runtime pairs $\{(n_i, t_i)\}$ against candidate models corresponding to the six classes: 
 constant, $\log n$, $n$, $n\log n$, $n^2$, and $n^3$. 
 Model parameters are estimated by least squares, and we select the best class using a penalized 
 objective that discourages unnecessarily high-order models under limited scales (Section~\ref{subsubsec:penalty}). 
 We additionally compute a fitting confidence score based on the margin between the best and second-best 
 penalized scores, which is used in the fusion correction rule. 
 
 \subsubsection{LLM Inference Settings} 
 \label{subsubsec:setup_llm} 
 For LLM-only and Fusion, we use a fixed prompt template with a constrained answer set (the six classes) 
 and require a structured output (JSON) to improve parsability. 
 We run inference with deterministic decoding (e.g., temperature set to 0) whenever supported. 
 Invalid or out-of-set outputs are mapped to \texttt{unknown} and treated as misclassifications in accuracy 
 and macro-F1 unless otherwise specified. 
 
 \subsubsection{Evaluation Metrics} 
 \label{subsubsec:setup_metrics} 
 We report overall accuracy and macro-F1 as primary metrics. 
 To better reflect the hierarchical nature of complexity labels, we optionally report hierarchy-aware 
 scores (e.g., HC/HC$_w$) following prior work. 
 We also report per-class performance and analyze confusion patterns between adjacent classes 
 (e.g., \texttt{nlogn} vs.\ \texttt{quadratic}, \texttt{quadratic} vs.\ \texttt{cubic}).

\subsection{Some Common Mistakes}\label{SCM}
\begin{itemize}
\item The word ``data'' is plural, not singular.
\item The subscript for the permeability of vacuum $\mu_{0}$, and other common scientific constants, is zero with subscript formatting, not a lowercase letter ``o''.
\item In American English, commas, semicolons, periods, question and exclamation marks are located within quotation marks only when a complete thought or name is cited, such as a title or full quotation. When quotation marks are used, instead of a bold or italic typeface, to highlight a word or phrase, punctuation should appear outside of the quotation marks. A parenthetical phrase or statement at the end of a sentence is punctuated outside of the closing parenthesis (like this). (A parenthetical sentence is punctuated within the parentheses.)
\item A graph within a graph is an ``inset'', not an ``insert''. The word alternatively is preferred to the word ``alternately'' (unless you really mean something that alternates).
\item Do not use the word ``essentially'' to mean ``approximately'' or ``effectively''.
\item In your paper title, if the words ``that uses'' can accurately replace the word ``using'', capitalize the ``u''; if not, keep using lower-cased.
\item Be aware of the different meanings of the homophones ``affect'' and ``effect'', ``complement'' and ``compliment'', ``discreet'' and ``discrete'', ``principal'' and ``principle''.
\item Do not confuse ``imply'' and ``infer''.
\item The prefix ``non'' is not a word; it should be joined to the word it modifies, usually without a hyphen.
\item There is no period after the ``et'' in the Latin abbreviation ``et al.''.
\item The abbreviation ``i.e.'' means ``that is'', and the abbreviation ``e.g.'' means ``for example''.
\end{itemize}
An excellent style manual for science writers is \cite{b7}.

\subsection{Authors and Affiliations}\label{AAA}
\textbf{The class file is designed for, but not limited to, six authors.} A 
minimum of one author is required for all conference articles. Author names 
should be listed starting from left to right and then moving down to the 
next line. This is the author sequence that will be used in future citations 
and by indexing services. Names should not be listed in columns nor group by 
affiliation. Please keep your affiliations as succinct as possible (for 
example, do not differentiate among departments of the same organization).

\subsection{Identify the Headings}\label{ITH}
Headings, or heads, are organizational devices that guide the reader through 
your paper. There are two types: component heads and text heads.

Component heads identify the different components of your paper and are not 
topically subordinate to each other. Examples include Acknowledgments and 
References and, for these, the correct style to use is ``Heading 5''. Use 
``figure caption'' for your Figure captions, and ``table head'' for your 
table title. Run-in heads, such as ``Abstract'', will require you to apply a 
style (in this case, italic) in addition to the style provided by the drop 
down menu to differentiate the head from the text.

Text heads organize the topics on a relational, hierarchical basis. For 
example, the paper title is the primary text head because all subsequent 
material relates and elaborates on this one topic. If there are two or more 
sub-topics, the next level head (uppercase Roman numerals) should be used 
and, conversely, if there are not at least two sub-topics, then no subheads 
should be introduced.

\subsection{Figures and Tables}\label{FAT}
\paragraph{Positioning Figures and Tables} Place figures and tables at the top and 
bottom of columns. Avoid placing them in the middle of columns. Large 
figures and tables may span across both columns. Figure captions should be 
below the figures; table heads should appear above the tables. Insert 
figures and tables after they are cited in the text. Use the abbreviation 
``Fig.~\ref{fig}'', even at the beginning of a sentence.

\begin{table}[htbp]
\caption{Table Type Styles}
\begin{center}
\begin{tabular}{|c|c|c|c|}
\hline
\textbf{Table}&\multicolumn{3}{|c|}{\textbf{Table Column Head}} \\
\cline{2-4} 
\textbf{Head} & \textbf{\textit{Table column subhead}}& \textbf{\textit{Subhead}}& \textbf{\textit{Subhead}} \\
\hline
copy& More table copy$^{\mathrm{a}}$& &  \\
\hline
\multicolumn{4}{l}{$^{\mathrm{a}}$Sample of a Table footnote.}
\end{tabular}
\label{tab1}
\end{center}
\end{table}

\begin{figure}[htbp]
\centerline{\includegraphics{fig1.png}}
\caption{Example of a figure caption.}
\label{fig}
\end{figure}

Figure Labels: Use 8 point Times New Roman for Figure labels. Use words 
rather than symbols or abbreviations when writing Figure axis labels to 
avoid confusing the reader. As an example, write the quantity 
``Magnetization'', or ``Magnetization, M'', not just ``M''. If including 
units in the label, present them within parentheses. Do not label axes only 
with units. In the example, write ``Magnetization (A/m)'' or ``Magnetization 
\{A[m(1)]\}'', not just ``A/m''. Do not label axes with a ratio of 
quantities and units. For example, write ``Temperature (K)'', not 
``Temperature/K''.

\section*{Acknowledgment}

The preferred spelling of the word ``acknowledgment'' in America is without 
an ``e'' after the ``g''. Avoid the stilted expression ``one of us (R. B. 
G.) thanks $\ldots$''. Instead, try ``R. B. G. thanks$\ldots$''. Put sponsor 
acknowledgments in the unnumbered footnote on the first page.

\section*{References}

Please number citations consecutively within brackets \cite{b1}. The 
sentence punctuation follows the bracket \cite{b2}. Refer simply to the reference 
number, as in \cite{b3}---do not use ``Ref. \cite{b3}'' or ``reference \cite{b3}'' except at 
the beginning of a sentence: ``Reference \cite{b3} was the first $\ldots$''

Number footnotes separately in superscripts. Place the actual footnote at 
the bottom of the column in which it was cited. Do not put footnotes in the 
abstract or reference list. Use letters for table footnotes.

Unless there are six authors or more give all authors' names; do not use 
``et al.''. Papers that have not been published, even if they have been 
submitted for publication, should be cited as ``unpublished'' \cite{b4}. Papers 
that have been accepted for publication should be cited as ``in press'' \cite{b5}. 
Capitalize only the first word in a paper title, except for proper nouns and 
element symbols.

For papers published in translation journals, please give the English 
citation first, followed by the original foreign-language citation \cite{b6}.

\begin{thebibliography}{00}
\bibitem{b1} G. Eason, B. Noble, and I. N. Sneddon, ``On certain integrals of Lipschitz-Hankel type involving products of Bessel functions,'' Phil. Trans. Roy. Soc. London, vol. A247, pp. 529--551, April 1955.
\bibitem{b2} J. Clerk Maxwell, A Treatise on Electricity and Magnetism, 3rd ed., vol. 2. Oxford: Clarendon, 1892, pp.68--73.
\bibitem{b3} I. S. Jacobs and C. P. Bean, ``Fine particles, thin films and exchange anisotropy,'' in Magnetism, vol. III, G. T. Rado and H. Suhl, Eds. New York: Academic, 1963, pp. 271--350.
\bibitem{b4} K. Elissa, ``Title of paper if known,'' unpublished.
\bibitem{b5} R. Nicole, ``Title of paper with only first word capitalized,'' J. Name Stand. Abbrev., in press.
\bibitem{b6} Y. Yorozu, M. Hirano, K. Oka, and Y. Tagawa, ``Electron spectroscopy studies on magneto-optical media and plastic substrate interface,'' IEEE Transl. J. Magn. Japan, vol. 2, pp. 740--741, August 1987 [Digests 9th Annual Conf. Magnetics Japan, p. 301, 1982].
\bibitem{b7} M. Young, The Technical Writer's Handbook. Mill Valley, CA: University Science, 1989.
\bibitem{b8} D. P. Kingma and M. Welling, ``Auto-encoding variational Bayes,'' 2013, arXiv:1312.6114. [Online]. Available: https://arxiv.org/abs/1312.6114
\bibitem{b9} S. Liu, ``Wi-Fi Energy Detection Testbed (12MTC),'' 2023, gitHub repository. [Online]. Available: https://github.com/liustone99/Wi-Fi-Energy-Detection-Testbed-12MTC
\bibitem{b10} ``Treatment episode data set: discharges (TEDS-D): concatenated, 2006 to 2009.'' U.S. Department of Health and Human Services, Substance Abuse and Mental Health Services Administration, Office of Applied Studies, August, 2013, DOI:10.3886/ICPSR30122.v2
\bibitem{b11} K. Eves and J. Valasek, ``Adaptive control for singularly perturbed systems examples,'' Code Ocean, Aug. 2023. [Online]. Available: https://codeocean.com/capsule/4989235/tree
\end{thebibliography}

\vspace{12pt}
\color{red}
IEEE conference templates contain guidance text for composing and formatting conference papers. Please ensure that all template text is removed from your conference paper prior to submission to the conference. Failure to remove the template text from your paper may result in your paper not being published.

\end{document}
